\section{The pro-Matlis target: Matlis principal parts and the continuous cotangent module}\label{sec:pro-matlis-target}\label{sec:app-matlis-principal-parts}

Configuration observables on the brane are the polynomial trace
classes \(O_f=\overline{\operatorname{Tr}f(\phi_1,\phi_2)}\): Taylor
in the matrix-valued normal coordinates, forming the PBW
special-fibre alphabet.  Their conjugate cotangent momenta, by
Grothendieck local-cohomology duality, are the residues
\(\rho_{a,b}=z_1^{-a-1}z_2^{-b-1}\,dz_1dz_2\), top local-cohomology
generators of \(R=\C[[z_1,z_2]]\) at the maximal ideal \((z_1,z_2)\).
The brane measures positions polynomially; its cotangent momenta live as
residues.

The \(\mathfrak m\)-adic continuous dual of the reduced Hamiltonian Lie
algebra \(\mathfrak h=R/\C\cdot1\) sits inside
\(H^2_{\mathfrak m}(R)\,dz_1dz_2\) as the scalar annihilator of the
constant line; the residue pairing identifies it canonically with the
direct sum \(\mc P=\bigoplus_{a+b>0}\C\,\rho_{a,b}\).  The Hamiltonian
coadjoint action, transported through this residue identification, is
the principal-part residue formula
\(z_1^p z_2^q\cdot\rho_{a,b}=-(pb-qa+p-q)\,\rho_{a-p+1,b-q+1}\), with
negative-index targets zero.  The polynomial PBW one-\(\psi\)
descendants \(\widetilde\Psi_{a,b}\) carry the same labels \((a,b)\) and
a strictly different \(\mathfrak h\)-module structure: linear
Hamiltonians act locally nilpotently on the polynomial side and never
locally nilpotently on \(\mc P\), so no relabelling of bases
interpolates.  Fourier--Rees conjugation supplies the only universal
label bridge, and only after a polarization change: it sends
\(\mathcal O_\hbar=\mathcal D_\hbar/(\partial_1,\partial_2)\) to
\(\Delta^{\mathsf F}_\hbar=\mathcal D_\hbar/(z_1,z_2)\), whose \v{C}ech
realization sits inside the standard local-cohomology lattice as
\(\bigoplus\hbar^{a+b}\C[[\hbar]]\rho_{a,b}\); the full lattice
\(\bigoplus\C[[\hbar]]\rho_{a,b}\) is recovered only after \(\hbar\) is
inverted.

Source conventions: Matlis duality \cite{matlis-injective},
Grothendieck--Hartshorne local cohomology
\cite{hartshorne-local-cohomology}, and the power-series residue model
of Kunz--Cox--Dickenstein, Chapter~5 \cite{kunz-residues-duality}.

\subsection{The defect polarization}
\label{ssec:app-defect-polarization}

A canonical-transformation BV theory on the formal symplectic polydisk
\((\widehat{\C^2}_0,dz_1\wedge dz_2)\) carries two coordinate types:
Hamiltonian-dual coordinates \(c^I\), which generate canonical motion,
and shifted-cotangent coordinates \(u_I\), which carry antifields.  The
brane line \(L=\R_t\times\{s=z_1=z_2=0\}\) reads only the second
alphabet, through the trace substitution of
Lemma~\ref{lem:formal-stalk-trace}: a holomorphic Hamiltonian
\(f\in\mathfrak h\) becomes a brane configuration observable
\(O_f=\overline{\operatorname{Tr}f(\phi_1,\phi_2)}\), a \emph{polynomial}
in the Chan--Paton-valued normal coordinates.  Configuration observables
are therefore Taylor functions of the Chan--Paton pair.  The cotangent
momentum dual to \(O_f\) is constrained by three properties: it pairs
with Hamiltonian observables to a scalar; it is continuous for the
\(\mathfrak m\)-adic topology on \(\mathfrak h\), where
\(\mathfrak m=(z_1,z_2)\); it intertwines the Hamiltonian coadjoint
action coming from the Poisson bracket on the polydisk.  Those three
constraints together select a polarization that is \emph{not} a copy of
polynomials: the cotangent direction is residue-valued.

\begin{thm}[Defect Polarization]
\label{thm:app-defect-polarization}
  On the formal symplectic polydisk \((\widehat{\C^2}_0,dz_1\wedge dz_2)\),
  the brane defect line \(L=\R_t\times\{s=z_1=z_2=0\}\) carries a
  canonical polarization with the following two sides.
  \begin{enumerate}
    \item \emph{Configuration side: polynomial / Taylor.}  Hamiltonian
      configuration observables on the brane are the polynomial trace
      family
      \[
        O_f=\overline{\operatorname{Tr}f(\phi_1,\phi_2)},\qquad
        f\in\mathfrak h=\C[[z_1,z_2]]/\C\cdot1 .
      \]
      The Lie algebra acting through these observables is the reduced
      polynomial Hamiltonian sub-Lie algebra
      \(\bar A=\C[z_1,z_2]/\C\cdot1\subset\mathfrak h\), each generator
      represented by a constant Hamiltonian vector field on the
      polynomial PBW special-fibre alphabet.  These observables form the
      PBW one-\(\psi\) shadow span:
      \(\widetilde\Psi_{a,b}\in\bar A_{\mathrm{desc}}\), bidegree
      \((a,b)\), with linear Hamiltonians acting locally nilpotently.
    \item \emph{Cotangent side: distributional / residual.}  Cotangent
      momenta on the same brane are the Grothendieck--Matlis principal
      parts
      \[
        \mc P=\bigoplus_{\substack{a,b\geq0\\a+b>0}}\C\,\rho_{a,b},
        \qquad
        \rho_{a,b}=z_1^{-a-1}z_2^{-b-1}\,dz_1dz_2,
      \]
      basis-paired with the Taylor monomial dual by the Grothendieck
      residue
      \(\langle\rho_{a,b},z_1^m z_2^n\rangle=\delta_{a,m}\delta_{b,n}\),
      and acted on by Hamiltonians through the principal-part coadjoint
      formula
      \[
        z_1^p z_2^q\cdot\rho_{a,b}
          =-(pb-qa+p-q)\rho_{a-p+1,b-q+1},
      \]
      with negative-index targets zero.  Linear Hamiltonians act on
      \(\mc P\) by raising pole order, never locally nilpotently.
  \end{enumerate}
  The two sides are not two bases of one \(\mathfrak h\)-module.  The
  bidegree relabelling
  \(\widetilde\Psi_{a,b}\leftrightarrow\rho_{a,b}\) carries the same
  index set \((a,b)\) but does not intertwine the two Hamiltonian
  actions: linear-Hamiltonian local nilpotence holds on the polynomial
  side and fails on the residual side, so any same-action module
  isomorphism is impossible
  (Theorem~\ref{thm:app-matlis-polynomial-descendant-obstruction}).  The
  Moyal coadjoint deformation acts on the residual side; the polynomial
  one-\(\psi\) shadow does not intertwine that action even after
  inverting \(\hbar\), and the Fourier--Rees label bridge of
  Theorem~\ref{thm:app-matlis-rees-fourier-bridge} is the unique
  universal interpolation, available only after a Fourier polarization
  change of the action itself.

  In particular, the continuous dual of the Hamiltonian Lie algebra of
  the defect is the residual side:
  \[
    \mathfrak h^\vee_{\mathrm{cont}}\cong\mc P
    \qquad\text{(Proposition~\ref{prop:app-matlis-realization})},
  \]
  and the residue-coadjoint pairing is the unique perfect continuous
  pairing intertwining the two sides
  (Theorem~\ref{thm:app-matlis-residue-rigidity}).
\end{thm}

\begin{proof}
  The configuration side is
  Lemma~\ref{lem:formal-stalk-trace}: the gauge-invariant content of
  the Chan--Paton pair on the brane is the trace family
  \(\{O_f\}_{f\in\bar A}\), with \(\bar A=\C[z_1,z_2]/\C\cdot1\) acting
  by Hamiltonian vector fields on the polynomial PBW special-fibre
  alphabet \(\widetilde\Psi_{a,b}\); local nilpotence of the linear
  Hamiltonians \(z_1,z_2\) on this side is
  Theorem~\ref{thm:app-matlis-polynomial-descendant-obstruction}.

  On the cotangent side, the basis identity
  \(\mc P=\bigoplus_{a+b>0}\C\,\rho_{a,b}\) and the residue Kronecker
  pairing
  \(\langle\rho_{a,b},z_1^m z_2^n\rangle=\delta_{a,m}\delta_{b,n}\) are
  Definition~\ref{def:app-matlis-principal-parts}.  The continuous-dual
  identification \(\mathfrak h^\vee_{\mathrm{cont}}\cong\mc P\) is
  Proposition~\ref{prop:app-matlis-realization}.  The principal-part
  coadjoint action with monomial coefficient \(-(pb-qa+p-q)\) is
  Proposition~\ref{prop:app-matlis-coadjoint-action} and
  Lemma~\ref{lem:app-matlis-coadjoint-continuity}; iteration on linear
  Hamiltonians produces the rising-factorial
  \((-1)^N(b+1)\cdots(b+N)\rho_{a,b+N}\) in
  Theorem~\ref{thm:app-matlis-polynomial-descendant-obstruction},
  excluding local nilpotence.

  Linear Hamiltonians act locally nilpotently on the polynomial side and
  not on the residual side, by the same theorem, so no same-action
  module isomorphism interpolates.  The unique universal label bridge
  between the two sides is Theorem~\ref{thm:app-matlis-rees-fourier-bridge};
  that the bridge requires a Fourier polarization change of action is
  the obstruction clause of the same theorem and
  Corollary~\ref{cor:app-matlis-no-flat-same-action-interpolation}.
  The residue-coadjoint pairing is unique up to one scalar by
  Theorem~\ref{thm:app-matlis-residue-rigidity}.
\end{proof}

\begin{rmk}[Computational verification of the polarization]
\label{rmk:app-defect-polarization-numerics}
  The coadjoint coefficient \(-(pb-qa+p-q)\) on the residual side is
  verified by \texttt{scripts/check\_one\_psi\_homology.py}, which
  computes \(\operatorname{ad}^\vee_{z_1^pz_2^q}(\delta_{a,b})\)
  symbolically on Taylor-monomial test labels, identifies the
  one-functional support, and confirms
  \(-(pb-qa+p-q)\rho_{a-p+1,b-q+1}\) on every quadruple \((p,q,a,b)\)
  with exponents \(\le 5\) and \(p+q,a+b>0\): \(1225\) cases, all match.
  The same script attacks the candidate label map
  \(\widetilde\Psi_{a,b}\mapsto\rho_{a,b}\) by the exact Moyal coadjoint
  action: \(1200\) of \(1225\) cases register a PBW-vs-Moyal mismatch,
  and \(958\) cases carry genuine higher Moyal coadjoint terms which the
  polynomial one-\(\psi\) action does not produce at any order in
  \(\hbar\).  The smallest explicit obstruction is the pair
  \[
    \operatorname{ad}^\vee_{z_1^4}(\delta_{1,1})
      =-24\,\hbar^2\,\delta_{0,4},
    \qquad
    \widetilde\Psi_{1,1}\mapsto4\,\widetilde\Psi_{4,0}
    \quad\text{(polynomial side)};
  \]
  the two functionals are supported at different bidegrees and at
  different powers of \(\hbar\).  The companion script
  \texttt{scripts/check\_moyal\_coefficients.py} verifies the Moyal
  star-bracket coefficients on monomials up to exponents \(\le 10\) and
  orders \(\le 11\), with the antisymmetric odd-order tower
  \(1,\ \tfrac1{24},\ \tfrac1{1920},\ \tfrac1{322560},\ \dots\) matching
  the closed-form \(\tfrac{1}{2^{r-1}\,r!}\) for odd \(r\).  Both
  scripts run from the command line with no external dependencies.
\end{rmk}

\subsection{Construction of the residual side}
\label{ssec:app-matlis-construction}

\begin{defn}[Topology and reduced Hamiltonians]
\label{def:app-matlis-topology}
  Put
  \[
    R=\C[[z_1,z_2]],\qquad \mathfrak m=(z_1,z_2),\qquad
    \mathfrak h=R/\C\cdot1 .
  \]
  The topology on \(R\) is the \(\mathfrak m\)-adic topology.  The
  topology on \(\mathfrak h\) is the quotient topology, equivalently
  the topology transported through the zero-constant section
  \[
    \mathfrak h\simeq \mathfrak m,\qquad
    \overline f\longmapsto f-f(0).
  \]
  A fundamental system of neighbourhoods of zero is
  \[
    F^N\mathfrak h=\mathfrak m^N,\qquad N\geq1,
  \]
  via this splitting.  The Lie bracket is the Poisson bracket
  \[
    \{f,g\}=\partial_{z_1}f\,\partial_{z_2}g
      -\partial_{z_2}f\,\partial_{z_1}g
  \]
  followed by projection modulo constants.  The adjoint maps are continuous
  for this topology: if \(f\in\mathfrak m\), then
  \(\{f,\mathfrak m^{N+1}\}\subset \mathfrak m^N\).  The constant class
  acts by zero and is not part of \(\mathfrak h\).
\end{defn}

\begin{lem}[Continuous dual]
\label{lem:app-matlis-continuous-dual}
  The continuous dual of \(\mathfrak h\), with \(\C\) discrete, is
  \[
    \mathfrak h^\vee_{\mathrm{cont}}
      =\bigcup_{N\geq1}\bigl(\mathfrak h/F^N\mathfrak h\bigr)^\vee
      =\bigoplus_{\substack{a,b\geq0\\a+b>0}}\C\,\delta_{a,b},
  \]
  where
  \[
    \delta_{a,b}(z_1^m z_2^n)=\delta_{a,m}\delta_{b,n}.
  \]
\end{lem}

\begin{proof}
  A linear functional on the \(\mathfrak m\)-adic vector space
  \(\mathfrak h\simeq\mathfrak m\) is continuous exactly when its
  kernel contains \(F^N\mathfrak h=\mathfrak m^N\) for some \(N\).
  Such a functional factors through the finite-dimensional quotient
  \(\mathfrak m/\mathfrak m^N\), hence is a finite linear combination
  of Taylor coefficient functionals.  Conversely every finite linear
  combination of the displayed \(\delta_{a,b}\) kills all sufficiently
  high powers of \(\mathfrak m\), hence is continuous.
\end{proof}

\begin{defn}[Reduced Matlis principal parts]
\label{def:app-matlis-principal-parts}
  The top local cohomology module of \(R\) is written in \v{C}ech
  Laurent form as
  \[
    H^2_{\mathfrak m}(R)\,dz_1dz_2
      =\bigoplus_{a,b\geq0}
        \C\,z_1^{-a-1}z_2^{-b-1}\,dz_1dz_2 .
  \]
  Equivalently this is the \v{C}ech quotient
  \[
    R[z_1^{-1},z_2^{-1}]/(R[z_1^{-1}]+R[z_2^{-1}])
  \]
  after multiplication by \(dz_1dz_2\); the displayed direct sum is
  the standard representative basis in which both exponents are
  strictly negative.
  The basis vector \(\rho_{0,0}=z_1^{-1}z_2^{-1}dz_1dz_2\) is not zero in
  \(H^2_{\mathfrak m}(R)\,dz_1dz_2\); it is the functional which evaluates
  the constant coefficient.  The reduced principal-part module is the
  \(\C\)-linear residue annihilator of constants,
  \[
    \mc P=\operatorname{Ann}_{\mathrm{res}}(\C\cdot1)
      =\bigoplus_{\substack{a,b\geq0\\a+b>0}}\C\,\rho_{a,b},
      \qquad
    \rho_{a,b}=z_1^{-a-1}z_2^{-b-1}\,dz_1dz_2 .
  \]
  Equip \(\mc P\) with the locally finite direct-sum topology
  \[
    \mc P=\varinjlim_N \mc P_{\leq N},\qquad
    \mc P_{\leq N}
      =\bigoplus_{\substack{a,b\geq0\\0<a+b\leq N}}\C\,\rho_{a,b},
  \]
  and give \(\mathfrak h^\vee_{\mathrm{cont}}\) the corresponding
  inductive-limit topology through the basis \(\delta_{a,b}\). Thus the
  residue identification below is a continuous map of the declared
  locally finite topological vector spaces, not a bare algebraic
  relabelling.  The space \(\mc P\) is a Hamiltonian-coadjoint module; it is
  not an \(R\)-submodule of \(H^2_{\mathfrak m}(R)\,dz_1dz_2\), since
  multiplication by \(z_i\) can hit the excluded scalar-residue line.
  The residue pairing is
  \[
    \langle \rho_{a,b},z_1^m z_2^n\rangle
      =\operatorname{Res}_{z_1=z_2=0}
        z_1^{m-a-1}z_2^{n-b-1}\,dz_1dz_2
      =\delta_{a,m}\delta_{b,n}.
  \]
\end{defn}

\begin{prop}[Matlis realization of the continuous dual]
\label{prop:app-matlis-realization}
  The residue pairing induces a continuous vector-space isomorphism
  \[
    \mc P\xrightarrow{\sim}\mathfrak h^\vee_{\mathrm{cont}},
    \qquad
    \rho_{a,b}\longmapsto\delta_{a,b}.
  \]
  It is the scalar-reduced restriction of the Matlis evaluation
  pairing for the complete regular local ring \(R\): after the
  trivialization \(\omega_R=R\,dz_1dz_2\),
  \(H^2_{\mathfrak m}(R)\,dz_1dz_2\cong E_R(\C)\).  The subspace
  \(\mc P=\operatorname{Ann}_{\mathrm{res}}(\C\cdot1)\) is the
  \(\C\)-linear scalar-annihilator dual to
  \(\mathfrak h=R/\C\cdot1\).  The \(R\)-module structure on
  \(E_R(\C)\) does not restrict to \(\mc P\); multiplication can hit the
  excluded scalar residue line.  The Hamiltonian coadjoint action is
  transported separately by the residue pairing
  (Proposition~\ref{prop:app-matlis-coadjoint-action}) and is the
  structure used by the brane cotangent sector.
\end{prop}

\begin{proof}
  The displayed residue formula identifies each basis vector
  \(\rho_{a,b}\), \(a+b>0\), with the continuous Taylor coefficient
  functional \(\delta_{a,b}\).  Lemma~\ref{lem:app-matlis-continuous-dual}
  shows that these coefficient functionals form the whole continuous
  dual.  The exclusion of \(\rho_{0,0}\) is exactly the condition that
  the functional vanish on the scalar line \(\C\cdot1\subset R\); this is
  the continuous dual of the quotient \(R/\C\cdot1\).  On finite stages
  the map sends \(\mc P_{\leq N}\) isomorphically to the functionals
  vanishing on \(\mathfrak m^{N+1}\), so both the map and its inverse are
  continuous for the declared locally finite topologies.
\end{proof}

\begin{prop}[Residue coadjoint action]
\label{prop:app-matlis-coadjoint-action}
  Under the isomorphism of
  Proposition~\ref{prop:app-matlis-realization}, the coadjoint action
  of \(f\in\mathfrak h\) is the principal-part residue action
  \[
    f\cdot\rho=\pi_{\mathrm{pp}}\{f,\rho\},
  \]
  where \(\pi_{\mathrm{pp}}\) projects Laurent two-forms to the span
  of the principal parts.  On monomials,
  \begin{equation}
  \label{eq:app-matlis-coadjoint-monomial}
    z_1^p z_2^q\cdot\rho_{a,b}
      =-(pb-qa+p-q)\,\rho_{a-p+1,b-q+1},
  \end{equation}
  with the convention that any term with a negative index is zero.  If the
  displayed target is \(\rho_{0,0}\), then the coefficient \(pb-qa+p-q\)
  is zero; hence the scalar residue line is not produced from \(\mc P\),
  and the action preserves the scalar-annihilator subspace.  The
  coefficient \(-(pb-qa+p-q)\) is the residue side of the defect
  polarization (Theorem~\ref{thm:app-defect-polarization}); on the
  polynomial PBW one-\(\psi\) shadow the analogous coefficient is the
  symmetric \(pb-qa\), and the two formulas differ already at the linear
  Hamiltonians.
\end{prop}

\begin{proof}
  Let \(X_f=(\partial_{z_1}f)\partial_{z_2}
  -(\partial_{z_2}f)\partial_{z_1}\).  Hamiltonian vector fields preserve
  the volume form \(dz_1dz_2\).  For a Laurent coefficient \(\phi\) and
  \(g\in R\), the residue of a total \(z_i\)-derivative is zero, hence
  \[
    0=\operatorname{Res}\bigl(L_{X_f}(\phi g\,dz_1dz_2)\bigr)
      =\operatorname{Res}\bigl((\{f,\phi\}g+\phi\{f,g\})\,dz_1dz_2\bigr).
  \]
  Therefore the residue pairing satisfies integration by parts,
  \[
    \langle \{f,\rho\},g\rangle
      =-\langle \rho,\{f,g\}\rangle ,
  \]
  which is exactly the coadjoint sign convention
  \((\operatorname{ad}^{\vee}_f\lambda)(g)=-\lambda(\{f,g\})\).

  For \(f=z_1^p z_2^q\) and
  \(\rho_{a,b}=z_1^{-a-1}z_2^{-b-1}\,dz_1dz_2\), expand the bracket.
  The two terms come from \(\partial_{z_1}f\cdot\partial_{z_2}\phi\)
  and \(-\partial_{z_2}f\cdot\partial_{z_1}\phi\):
  \[
  \begin{aligned}
    \{f,z_1^{-a-1}z_2^{-b-1}\}
      &=
      pz_1^{p-1}z_2^q\,(-(b+1))z_1^{-a-1}z_2^{-b-2}  \\
      &\quad
      -qz_1^p z_2^{q-1}\,(-(a+1))z_1^{-a-2}z_2^{-b-1} \\
      &=\bigl(-p(b+1)+q(a+1)\bigr)z_1^{p-a-2}z_2^{q-b-2} \\
      &=-(pb-qa+p-q)z_1^{p-a-2}z_2^{q-b-2}.
  \end{aligned}
  \]
  Multiplication by \(dz_1dz_2\) gives
  \(\rho_{a-p+1,b-q+1}\) precisely when the two new indices are
  nonnegative; otherwise the principal-part projection gives zero.  If
  \(a-p+1=b-q+1=0\), then \(a=p-1\) and \(b=q-1\), and the coefficient
  \(pb-qa+p-q=p(q-1)-q(p-1)+p-q\) reduces to zero by direct expansion.
  Thus the action preserves the scalar-annihilator \(\mc P\).  The
  signed coefficient is verified by direct computation against
  \texttt{check\_one\_psi\_homology.py} on \(1225\) cases with exponents
  \(\le 5\).
\end{proof}

\begin{lem}[Continuity of the coadjoint action]
\label{lem:app-matlis-coadjoint-continuity}
  The formula of Proposition~\ref{prop:app-matlis-coadjoint-action}
  extends from polynomial Hamiltonians to all of
  \(\mathfrak h=\C[[z_1,z_2]]/\C\cdot1\).  The bilinear action
  \[
    \mathfrak h\times\mc P\longrightarrow\mc P
  \]
  is continuous for the \(\mathfrak m\)-adic topology on \(\mathfrak h\)
  and the locally finite direct-sum topology on \(\mc P\).
\end{lem}

\begin{proof}
  Fix \(\rho_{a,b}\).  In the monomial formula a term
  \(z_1^p z_2^q\cdot\rho_{a,b}\) can be nonzero only if
  \(p\leq a+1\) and \(q\leq b+1\).  Thus an arbitrary formal Hamiltonian
  acts on \(\rho_{a,b}\) through a finite jet, so the action is
  well-defined and continuous in the Hamiltonian variable.  Conversely,
  if \(a+b\leq N\), every nonzero target has total index
  \[
    (a-p+1)+(b-q+1)=a+b-(p+q)+2\leq N+1
  \]
  because \(p+q>0\) for a reduced Hamiltonian.  Hence each fixed
  Hamiltonian sends \(\mc P_{\leq N}\) into \(\mc P_{\leq N+1}\).  This is
  exactly continuity for the locally finite direct limit.
\end{proof}

\begin{lem}[Coadjoint centralizer]
\label{lem:app-matlis-coadjoint-centralizer}
  Every \(\mathfrak h\)-module endomorphism of \(\mc P\), with the
  coadjoint action above, is multiplication by one scalar.
\end{lem}

\begin{proof}
  The vector \(\rho_{1,0}\) is cyclic.  Indeed
  \(z_1^2\cdot\rho_{1,0}=-2\rho_{0,1}\); repeated action by the linear
  Hamiltonians \(z_1\) and \(z_2\) then reaches every
  \(\rho_{a,b}\), \(a+b>0\), with nonzero rising-factorial coefficient.

  Let \(T\) be an \(\mathfrak h\)-module endomorphism.  Since
  \(z_1z_2\cdot\rho_{a,b}=(a-b)\rho_{a,b}\), the vector
  \(T\rho_{1,0}\) lies in the eigenvalue-one subspace
  \[
    \bigoplus_{r\geq0}\C\,\rho_{r+1,r}.
  \]
  It is a finite sum because \(\mc P\) is a direct sum:
  \(T\rho_{1,0}=\sum_{r=0}^R c_r\rho_{r+1,r}\).  The Hamiltonian
  \(z_2^2\) kills \(\rho_{1,0}\), hence it kills \(T\rho_{1,0}\).  But
  for \(r\geq1\),
  \[
    z_2^2\cdot\rho_{r+1,r}
      =(2r+4)\rho_{r+2,r-1},
  \]
  while the \(r=0\) target has negative second index and is zero.  The
  displayed nonzero targets are linearly independent, so \(c_r=0\) for
  all \(r\geq1\).  Therefore \(T\rho_{1,0}=c_0\rho_{1,0}\), and cyclicity
  gives \(T=c_0\,\mathrm{id}_{\mc P}\).
\end{proof}

\subsection{Rigidity of the residue pairing}
\label{ssec:app-matlis-rigidity}

\begin{thm}[Residue pairing rigidity]
\label{thm:app-matlis-residue-rigidity}
  Up to multiplication by one scalar in \(\C^\times\), the residue
  pairing is the unique perfect continuous pairing
  \[
    \mc P\times\mathfrak h\longrightarrow\C
  \]
  which is equivariant for the coadjoint action above.  Here
  \emph{perfect} means that the adjoint map
  \(\mc P\to\mathfrak h^\vee_{\mathrm{cont}}\) is a topological
  vector-space isomorphism.  In particular the same uniqueness holds in
  the narrower class of total-degree-zero pairings.  Fixing the scalar by
  \[
    \langle\rho_{a,b},z_1^m z_2^n\rangle
      =\delta_{a,m}\delta_{b,n}
  \]
  is therefore the universal dual-module normalization for the
  principal-part cotangent sector.
\end{thm}

\begin{proof}
  Let \(\beta\) be another perfect continuous equivariant pairing.  It
  gives a topological vector-space isomorphism
  \[
    \Phi_\beta:\mc P\longrightarrow\mathfrak h^\vee_{\mathrm{cont}},
    \qquad
    \Phi_\beta(\rho)(g)=\beta(\rho,g).
  \]
  Equivariance of \(\beta\) says
  \[
    \beta(f\cdot\rho,g)=-\beta(\rho,\{f,g\}),
  \]
  so \(\Phi_\beta\) is an \(\mathfrak h\)-module map to the coadjoint
  module.  Let \(\Phi_{\mathrm{res}}\) be the residue isomorphism of
  Proposition~\ref{prop:app-matlis-realization}.  Then
  \(T=\Phi_{\mathrm{res}}^{-1}\Phi_\beta\) is an \(\mathfrak h\)-module
  automorphism of \(\mc P\).  By
  Lemma~\ref{lem:app-matlis-coadjoint-centralizer},
  \(T=c\,\mathrm{id}_{\mc P}\).  Since \(\Phi_\beta\) is an isomorphism,
  \(c\neq0\), and
  \(\beta=c\,\langle\,\cdot\,,\,\cdot\,\rangle_{\mathrm{res}}\).  The
  displayed Kronecker normalization sets \(c=1\).
\end{proof}

\subsection{Local-nilpotence obstruction to polynomial realization}
\label{ssec:app-matlis-polynomial-obstruction}

\begin{thm}[Polynomial descendant obstruction for principal parts]
\label{thm:app-matlis-polynomial-descendant-obstruction}
  Let \(M\) be a polynomial boundary descendant module on which the
  linear Hamiltonians \(z_1,z_2\) act locally nilpotently.  Then
  \(\mc P\) is not isomorphic to any nonzero submodule of \(M\) as an
  \(\mathfrak h\)-module.  In particular the polynomial PBW
  one-\(\psi\) descendants \(\widetilde\Psi_{a,b}\) cannot be identified
  with the principal parts \(\rho_{a,b}\) by an ordinary
  \(\mathfrak h\)-module isomorphism.
\end{thm}

\begin{proof}
  In any ordinary polynomial boundary model, the linear Hamiltonians
  act as constant coefficient Hamiltonian vector fields; iterating
  either \(z_1\) or \(z_2\) eventually exhausts the polynomial degree.
  Thus these actions are locally nilpotent: for every
  \(m\in M\) and every \(i\in\{1,2\}\) there exists \(N\geq1\) with
  \((z_i\cdot)^N m=0\).

  On \(\mc P\) the situation is rigorously the opposite.
  Proposition~\ref{prop:app-matlis-coadjoint-action} specialises to
  \[
    z_1\cdot\rho_{a,b}=-(b+1)\,\rho_{a,b+1},\qquad
    z_2\cdot\rho_{a,b}=(a+1)\,\rho_{a+1,b},
  \]
  and iteration produces the rising-factorial coefficient
  \[
    z_1^N\cdot\rho_{a,b}
      =(-1)^N(b+1)(b+2)\cdots(b+N)\,\rho_{a,b+N}\neq0
  \]
  for every \(N\geq0\); the same holds for \(z_2\).  Local nilpotence is
  preserved by module isomorphism and by restriction to submodules.
  Therefore no polynomial descendant module with locally nilpotent
  linear Hamiltonians realizes \(\mc P\), nor any nonzero submodule of
  \(\mc P\).  The obstruction is independent of \(\hbar\) and survives
  every label relabelling between the two alphabets.

  For the primitive PBW labels \(\widetilde\Psi_{a,b}\) the obstruction
  is already visible at the boundary of the index quadrant.  After the
  scalar descendant is removed, the special-fibre action of the linear
  Hamiltonians on the polynomial one-\(\psi\) shadow is
  \[
    z_1\cdot\widetilde\Psi_{a,b}=b\,\widetilde\Psi_{a,b-1},
    \qquad
    z_2\cdot\widetilde\Psi_{a,b}=-a\,\widetilde\Psi_{a-1,b}.
  \]
  Thus \(z_1\) kills \(\widetilde\Psi_{a,0}\) for every \(a>0\), whereas
  on the residual side \(z_1\cdot\rho_{a,0}=-\rho_{a,1}\neq0\).  No
  rescaling \(\widetilde\Psi_{a,b}\mapsto\lambda_{a,b}\rho_{a,b}\) with
  \(\lambda_{a,b}\neq0\) intertwines the two actions: the linear
  Hamiltonian \(z_1\) lowers the second index on the polynomial side
  and raises it on the residual side.
\end{proof}

\subsection{Universal Fourier--Rees label bridge}
\label{ssec:app-matlis-fourier-rees}

\begin{thm}[Universal Fourier--Rees construction and exact obstruction]
\label{thm:app-matlis-rees-fourier-bridge}
  Let
  \[
    \mathcal D_\hbar
    =
    \C[[\hbar]]
      \langle z_1,z_2,\partial_1,\partial_2\rangle
      \big/
      \bigl([\partial_i,z_j]=\hbar\delta_{ij},
      [z_i,z_j]=[\partial_i,\partial_j]=0\bigr)
  \]
  be the Rees Weyl algebra.  Let
  \[
    \mathcal O_\hbar=\C[z_1,z_2][[\hbar]]
      \cong\mathcal D_\hbar/(\partial_1,\partial_2)
  \]
  be the polynomial Rees module, and let
  \[
    \Delta^{\mathsf F}_\hbar=
      \mathcal D_\hbar/(z_1,z_2)
  \]
  be the Fourier delta Rees module, generated by a vector \(e_0\)
  killed by \(z_1,z_2\).  These quotients are left cyclic
  \(\mathcal D_\hbar\)-modules, with the indicated left ideals
  generated by \(\partial_1,\partial_2\) and by \(z_1,z_2\). PBW normal
  form gives \(\mathcal O_\hbar\) the free \(\C[[\hbar]]\)-basis
  \(z_1^az_2^b\) and gives \(\Delta^{\mathsf F}_\hbar\) the free basis
  \(\partial_1^a\partial_2^b e_0\); in particular both Rees modules are
  \(\C[[\hbar]]\)-flat.  The \v{C}ech local-cohomology lattice
  \[
    \Delta^{\mathrm{std}}_\hbar=
    H^2_{\mathfrak m}(\C[[z_1,z_2]])\,dz_1dz_2[[\hbar]]
  \]
  carries the standard action in which \(z_i\) multiplies and
  \(\partial_i=\hbar\partial_{z_i}\).  Sending
  \(e_0\mapsto\rho_{0,0}=z_1^{-1}z_2^{-1}dz_1dz_2\) realizes
  \(\Delta^{\mathsf F}_\hbar\) as the Rees sublattice
  \[
    \iota_{\mathsf F}(\Delta^{\mathsf F}_\hbar)
    =
    \bigoplus_{a,b\geq0}
      \C[[\hbar]]\,
      (-1)^{a+b}\hbar^{a+b}a!b!\,\rho_{a,b}
    \subset
    \Delta^{\mathrm{std}}_\hbar .
  \]
  This inclusion becomes an isomorphism after tensoring with
  \(\C((\hbar))\), but it is not an isomorphism of the two
  \(\C[[\hbar]]\)-lattices.  The Rees Fourier
  automorphism
  \[
    \mathsf F(z_i)=\partial_i,\qquad
    \mathsf F(\partial_i)=-z_i
  \]
  identifies the Fourier twist \(\mathsf F(\mathcal O_\hbar)\) with
  \(\Delta^{\mathsf F}_\hbar\).  This is the flat positive bridge: the
  action of \(P\in\mathcal D_\hbar\) on the polynomial source is carried
  to the action of \(\mathsf F(P)\) on the Fourier delta target.  On basis
  labels in the \v{C}ech
  realization,
  \[
    z_1^az_2^b
    \longmapsto
    \partial_1^a\partial_2^b e_0
    \overset{\iota_{\mathsf F}}{\longmapsto}
    \partial_1^a\partial_2^b\rho_{0,0}
    =
    (-1)^{a+b}\hbar^{a+b}a!b!\,\rho_{a,b}.
  \]
  After inverting \(\hbar\), this gives a filtered Fourier--Rees
  identification of label spaces between nonconstant polynomial labels and
  the reduced principal parts \(\rho_{a,b}\), \(a+b>0\), but only after the
  Fourier twist of action just stated.  Over the Rees lattice before
  inverting \(\hbar\), the construction records the degree-\(a+b\) shift
  by the explicit factor \(\hbar^{a+b}\).  The special fibre of the
  abstract Rees module keeps the labels
  \(\partial_1^a\partial_2^be_0\); the \v{C}ech inclusion into
  \(\Delta^{\mathrm{std}}_\hbar\) sends the positive-degree labels into
  \(\hbar\Delta^{\mathrm{std}}_\hbar\).  Thus the comparison is an
  associated-graded label match, not an isomorphism of ordinary
  local-cohomology lattices.

  The obstruction for the original Hamiltonian action is exact.  The
  Hamiltonian action on polynomial descendants is by vector fields
  \(X_f=\{f,-\}\) on \(\mathcal O_\hbar((\hbar))\).  In the Rees Weyl
  algebra the normalized linear operators are
  \(\hbar X_{z_1}=\partial_2\) and
  \(\hbar X_{z_2}=-\partial_1\).  Their Fourier transforms are
  \(\mathsf F(\hbar X_f)\), whereas the Matlis cotangent action on
  \(\mc P((\hbar))\) is still \(\hbar X_f\), acting on Laurent
  principal parts and projected to the reduced principal-part span.
  Already for the linear Hamiltonians,
  \[
    \hbar X_{z_1}=\partial_2,\qquad
    \hbar X_{z_2}=-\partial_1,
  \]
  but
  \[
    \mathsf F(\hbar X_{z_1})=-z_2,\qquad
    \mathsf F(\hbar X_{z_2})=z_1 .
  \]
  The transformed linear operators are locally nilpotent on
  \(\Delta^{\mathsf F}_\hbar\); the Rees-normalized coadjoint operators
  \[
    \hbar X_{z_1}\rho_{a,b}=-(b+1)\hbar\rho_{a,b+1},\qquad
    \hbar X_{z_2}\rho_{a,b}=(a+1)\hbar\rho_{a+1,b}
  \]
  are not locally nilpotent on any nonzero reduced principal part after
  inverting \(\hbar\).  Dividing by \(\hbar\) gives the same conclusion
  for the unnormalized Hamiltonian vector fields.  The
  Fourier--Rees construction therefore identifies the polynomial
  descendant module with \(\mc P\) as a module over a Fourier-conjugate
  copy of the reduced Hamiltonian Lie algebra; identification over the
  unconjugated action carries the same-action obstruction recorded
  below.  Any universal bridge selects one of two routes:
  Fourier-conjugate the action as above, or carry the same-action
  obstruction in the original Matlis coadjoint frame.
\end{thm}

\begin{proof}
  The displayed formulas for \(\mathsf F\) preserve the Rees Weyl
  relations, since
  \[
    [\mathsf F(\partial_i),\mathsf F(z_j)]
      =[-z_i,\partial_j]=\hbar\delta_{ij}.
  \]
  The polynomial module \(\mathcal O_\hbar\) is cyclic with generator
  \(1\) annihilated by \(\partial_1,\partial_2\).  After Fourier twist
  the cyclic generator is annihilated by \(z_1,z_2\), so the twisted
  module is \(\mathcal D_\hbar/(z_1,z_2)=\Delta^{\mathsf F}_\hbar\).
  In the \v{C}ech model, \(z_i\rho_{0,0}=0\), and applying
  \(\partial_i=\hbar\partial_{z_i}\) to \(\rho_{0,0}\) generates the
  Rees lattice displayed in the theorem.  It generates the full
  standard local-cohomology lattice only after \(\hbar\) is inverted.
  Therefore
  \(\mathsf F(\mathcal O_\hbar)\cong\Delta^{\mathsf F}_\hbar\), with the
  \v{C}ech realization \(\iota_{\mathsf F}\).

  Since \(\partial_i\) acts as \(\hbar\partial_{z_i}\) on Laurent
  representatives,
  \[
    \partial_1^a\partial_2^b\rho_{0,0}
    =
    \hbar^{a+b}
    \partial_{z_1}^a\partial_{z_2}^b
      (z_1^{-1}z_2^{-1}dz_1dz_2)
    =
    (-1)^{a+b}\hbar^{a+b}a!b!\rho_{a,b}.
  \]
  Removing the scalar label \(a=b=0\) gives the reduced label comparison
  after the stated Rees normalization.  Since the factors \(a!b!\) and
  \((-1)^{a+b}\) are units over \(\C\), the only genuine Rees-lattice
  shift is the power \(\hbar^{a+b}\).

  For the obstruction, compare the two linear actions.  On polynomial
  descendants the Rees-normalized operators
  \(\hbar X_{z_1}=\partial_2\) and
  \(\hbar X_{z_2}=-\partial_1\) are locally nilpotent; after inverting
  \(\hbar\), the same is true for \(X_{z_1}\) and \(X_{z_2}\).  Fourier
  sends the Rees-normalized operators to multiplication by \(-z_2\) and
  \(z_1\) on \(\Delta^{\mathsf F}_\hbar\); multiplication by \(z_i\)
  lowers pole order and is again locally nilpotent.  The coadjoint
  Matlis action is not the Fourier-transformed action.  It keeps the
  same vector fields, now acting on Laurent principal parts:
  \[
    \hbar X_{z_1}\rho_{a,b}
      =\hbar\partial_{z_2}\rho_{a,b}
      =-(b+1)\hbar\rho_{a,b+1},\qquad
    \hbar X_{z_2}\rho_{a,b}
      =-\hbar\partial_{z_1}\rho_{a,b}
      =(a+1)\hbar\rho_{a+1,b}.
  \]
  After tensoring with \(\C((\hbar))\), no nonzero vector in
  \(\mc P((\hbar))\) is killed by a high power of either linear
  Hamiltonian: for a finite nonzero linear combination, every power of
  \(X_{z_1}\) shifts each nonzero basis summand to a distinct
  \(\rho_{a,b+N}\), with nonzero rising-factorial coefficient, and the
  same argument applies to \(X_{z_2}\).  Local nilpotence is invariant
  under module isomorphism.
  Thus any Fourier--Rees comparison that preserves the original
  Hamiltonian action is impossible; the positive construction
  changes the polarization and is a flat \( \mathcal D_\hbar \)-module
  bridge only for the Fourier-conjugated action stated above.
\end{proof}

\begin{cor}[Same-action obstruction criterion]
\label{cor:app-matlis-no-flat-same-action-interpolation}
  Let \(M_\hbar\) be a \(\C[[\hbar]]\)-flat family whose generic fibre
  carries the polynomial descendant action of the linear Hamiltonians
  \(z_1,z_2\), so that \(X_{z_1}\) and \(X_{z_2}\) act locally
  nilpotently after tensoring with \(\C((\hbar))\).  Then
  \(M_\hbar((\hbar))\) is not isomorphic, as a module for the same
  Hamiltonian Lie algebra, to \(\mc P((\hbar))\) with the Matlis
  coadjoint action.  The same statement holds for the Rees-normalized
  operators \(\hbar X_{z_i}\), since \(\hbar\) is invertible on the
  generic fibre.  Equivalently, local nilpotence of the two linear
  Hamiltonians on the generic fibre is a necessary condition for a flat
  same-action Rees bridge, and \(\mc P((\hbar))\) violates that condition
  on every nonzero vector.  Any same-action comparison must therefore
  either change the Hamiltonian action by a Fourier twist, enlarge the
  boundary target to a distributional principal-part sector, or state a
  different deformation problem.
\end{cor}

\begin{proof}
  This is the local-nilpotence obstruction of
  Theorem~\ref{thm:app-matlis-polynomial-descendant-obstruction} after extension
  of scalars to \(\C((\hbar))\).  The two linear Hamiltonians remain
  locally nilpotent on the polynomial descendant generic fibre and
  remain non-locally-nilpotent on every nonzero element of
  \(\mc P((\hbar))\).  Multiplication by the invertible scalar
  \(\hbar\) cannot change local nilpotence.
\end{proof}

\begin{rmk}[Lattice content of the Fourier--Rees bridge]
\label{rmk:app-matlis-rees-bridge-scope}
  The shared index set \((a,b)\) is the Fourier--Rees
  associated-graded convention, with explicit factorial and
  \(\hbar^{a+b}\) normalization.  The Fourier twist gives
  \(\mathcal D_\hbar/(z_1,z_2)\), whose \v{C}ech realization is the Rees
  sublattice \(\bigoplus\hbar^{a+b}\C[[\hbar]]\rho_{a,b}\), not the full
  \(\bigoplus\C[[\hbar]]\rho_{a,b}\) before \(\hbar\) is inverted.
  Fourier sends the Hamiltonian vector fields to Fourier-transformed
  operators on \(\Delta^{\mathsf F}_\hbar\); the Matlis cotangent module
  uses the original coadjoint vector fields on principal parts.
\end{rmk}

